\documentclass[11pt,twocolumn,a4paper]{article}

\usepackage[margin=2.2cm]{geometry}
\usepackage{amsmath,amssymb,amsthm,mathtools}
\usepackage{bm}
\usepackage{graphicx}
\usepackage{hyperref}
\usepackage{natbib}
\usepackage{booktabs}
\usepackage{array}
\usepackage{enumitem}
\usepackage{float}
\usepackage{caption}
\captionsetup{font=small,labelfont=bf}

\usepackage[utf8]{inputenc}
\usepackage[T1]{fontenc}
\usepackage{lmodern}
\usepackage{microtype}

\newtheorem{theorem}{Theorem}[section]
\newtheorem{lemma}[theorem]{Lemma}
\newtheorem{proposition}[theorem]{Proposition}
\newtheorem{corollary}[theorem]{Corollary}
\theoremstyle{definition}
\newtheorem{definition}[theorem]{Definition}
\newtheorem{axiom}{Axiom}
\theoremstyle{remark}
\newtheorem{remark}[theorem]{Remark}

\hypersetup{colorlinks=true,linkcolor=blue,citecolor=blue,urlcolor=blue}

\newcommand{\Sfield}{S_{\mathrm{t}}}
\newcommand{\RA}{\alpha}
\newcommand{\Dec}{\delta}
\newcommand{\sky}{\mathsf{sky}}
\newcommand{\obs}{\mathrm{obs}}
\newcommand{\src}{\mathrm{src}}

\title{\textbf{The Emergent Celestial Light Field: \\
Astronomical Observation as Scalar-Field Measurement \\
on the Celestial Sphere}}

\author{Kundai Farai Sachikonye \\
\texttt{kundai.sachikonye@bitspark.com}}

\date{\today}

\begin{document}
\maketitle

\begin{abstract}
Classical astronomy treats the sky as a catalogue of objects and
assigns to each a flux, a position, a spectrum, and a set of
time-variability parameters. We reformulate the measurement as the
estimation of a single three-index scalar field
$\Sfield(\RA,\Dec,\omega)$ on the celestial sphere crossed with
temporal frequency. Every classical observable --- broadband
photometry, chromatographic spectroscopy (when time-resolved), stellar
occultation, pulsar timing, exoplanet transit depth, gravitational-wave
ringdown, satellite streak detection --- is a linear projection of this
field through a class-specific kernel. The measurement primitive is
subtraction of a known reference sky model from a time-series at a
fixed celestial address, followed by temporal Fourier transformation.
We prove four structural theorems: the observation primitive is linear
and shift-equivariant (Thm.~\ref{thm:linear}); any classical observable
of finite temporal support is recovered by kernel projection to within
the reference-model noise floor (Thm.~\ref{thm:projection}); object
identification is equivalent to cross-correlation of the pixel spectrum
against a class fingerprint (Thm.~\ref{thm:classification}); and
day--night distinction is dissolved at the field level because the
residual after subtraction is reference-model-invariant
(Thm.~\ref{thm:invariance}). Numerical validation across a synthetic
sky with 216 sources of five classes achieves
96\% class-recall at a false-positive rate below $5\times 10^{-3}$.
A browser-native implementation processes $96 \times 48$ pixels with
$256$ temporal samples at 15~Hz on a consumer laptop. The instrument
that measures $\Sfield$ replaces every optical telescope, spectrograph,
occultation pipeline, and photometric survey with a single endpoint.
\medskip

\noindent\textbf{Keywords:} scalar-field observation, differential
photometry, temporal spectrum, sky model, object identification,
classification by fingerprint, celestial sphere, shift-invariant
observation
\end{abstract}

%==============================================================================
\section{Introduction}
\label{sec:introduction}
%==============================================================================

\subsection{The object-flux ontology of classical astronomy}

Modern astronomical practice is built on a specific ontology: the sky
is a set of discrete objects, each characterised by a position, a
spectrum, a brightness history, and ancillary parameters such as
proper motion, radial velocity, polarisation, and classification
label \citep{binney2008}. Telescopes measure the objects; spectrographs
decompose their light; time-domain surveys record their variability;
occultation pipelines exploit the geometry of foreground bodies to
resolve angular structure. Each instrument class develops its own
pipeline, data format, and calibration methodology.

This ontology is pragmatic but not forced. Nothing in the physics of
electromagnetic radiation selects objects as primary. The quantity a
ground-based photon detector records is, at each instant and at each
pixel of its focal plane, a scalar: incident photon flux in a given
band. The object-as-primitive picture is a post-hoc reconstruction
from that scalar.

\subsection{An alternative: the field is the primitive}

We propose a different primitive. The sky is a scalar field on the
celestial sphere crossed with time: $I(\RA, \Dec, t)$. Subtract a
known reference sky model from $I$, and the residual's temporal
Fourier transform is a three-index scalar field
$\Sfield(\RA, \Dec, \omega)$ on $\mathbb{S}^2 \times \mathbb{R}_+$. We
call this the \emph{emergent light field}. Every classical observable
is a linear functional of $\Sfield$: the framework does not discard
objects, it derives them as emergent features of the field.

\subsection{What the reformulation gains}

Three practical gains follow from this change of variables.

\begin{itemize}[nosep]
\item \textbf{Observation primitive collapses to subtraction.}
Day-time, moonlit, and twilight observations become structurally
identical to dark-sky observations once the reference model is
subtracted. The Sun's scattered brightness enters the measurement as
an additive term known to the modeller, which vanishes in the residual.
\item \textbf{Object identification is a kernel cross-correlation, not
a detection threshold.} The spectral fingerprint of a pulsar, an
exoplanet-transit host, or a satellite streak is a known function of
$\omega$. Objects of each class are identified by the inner product
of the pixel spectrum against the class fingerprint.
\item \textbf{A single instrument replaces many.} Every classical
measurement is a projection of the field. The software pipeline
becomes: (i) compute $\Sfield$, (ii) apply the projection appropriate
to the query.
\end{itemize}

\subsection{Contributions}

\begin{enumerate}[nosep]
\item We formulate the emergent light field rigorously
(Section~\ref{sec:formulation}), define the subtraction observation
primitive, and establish its linearity and shift-equivariance
(Thm.~\ref{thm:linear}).
\item We prove that any classical observable with finite temporal
support is recovered as a linear projection of the field
(Thm.~\ref{thm:projection}), and that object identification by
spectral fingerprint is equivalent to kernel cross-correlation
(Thm.~\ref{thm:classification}).
\item We establish day--night equivalence at the field level: the
residual is invariant to the choice of reference model as long as the
model is accurate (Thm.~\ref{thm:invariance}).
\item We validate the framework numerically on a synthetic sky with
216 sources spanning five classes, achieving 96\% recall at false
positive rate below $5\times 10^{-3}$
(Section~\ref{sec:validation}).
\item We implement the instrument as a browser-native application
running at 15~Hz on consumer hardware
(Section~\ref{sec:implementation}).
\end{enumerate}

%==============================================================================
\section{Formulation}
\label{sec:formulation}
%==============================================================================

\subsection{The measured brightness time-series}

\begin{definition}[Measured brightness]
Let $(\RA, \Dec)$ index angular position on the celestial sphere
$\mathbb{S}^2$ and let $t \in [0, T]$ index time. The measured
brightness at a given pixel is
\begin{equation}
\label{eq:measurement}
I(\RA, \Dec, t) = I_\src(\RA, \Dec, t) + I_\sky(\RA, \Dec, t) + \varepsilon(\RA, \Dec, t),
\end{equation}
where $I_\src$ is the contribution from all celestial sources,
$I_\sky$ is the contribution from modellable atmospheric and
scattered-sunlight effects, and $\varepsilon$ is a zero-mean noise
term (photon shot noise plus instrument read noise).
\end{definition}

\subsection{The reference sky model}

\begin{definition}[Reference sky model]
A reference sky model $\widehat{I}_\sky(\RA, \Dec, t)$ is a
deterministic function encoding the modeller's best estimate of
$I_\sky$. Standard inputs: sun position and altitude, observer
latitude and longitude, time of day, aerosol optical depth, airglow
emission lines. Parametric forms include Rayleigh + Mie + zodiacal
radiative transfer \citep{noll2012}; variants include Krisciunas--Schaefer
moon models \citep{krisciunas1991} and empirical airglow models
\citep{patat2008}.
\end{definition}

\subsection{The residual and its spectrum}

\begin{definition}[Residual signal]
The residual is $R(\RA, \Dec, t) = I(\RA, \Dec, t) - \widehat{I}_\sky(\RA, \Dec, t)$.
\end{definition}

\begin{definition}[Emergent light field]
The emergent light field is the temporal power spectrum of the
residual:
\begin{equation}
\label{eq:field}
\Sfield(\RA, \Dec, \omega) = \left| \mathcal{F}_t [R(\RA, \Dec, \cdot)] (\omega) \right|^2,
\end{equation}
where $\mathcal{F}_t$ denotes the temporal Fourier transform. For a
time series of length $N$ sampled at interval $\Delta t$, $\omega$
ranges over $\{0, 1/(N\Delta t), 2/(N\Delta t), \ldots, 1/(2\Delta t)\}$
with Nyquist at $1/(2\Delta t)$.
\end{definition}

\subsection{Observation as linear map}

\begin{theorem}[Linearity and shift-equivariance]
\label{thm:linear}
The map from measured brightness to emergent light field, computed as
the power spectrum of the residual, is a bilinear form in the source
spectrum and shift-equivariant in time within each pixel.
\end{theorem}

\begin{proof}
Fix a pixel. The residual is a linear function of $I_\src$ (the
reference model cancels the $I_\sky$ contribution exactly if the model
is accurate; if not, a bias term $I_\sky - \widehat{I}_\sky$ is added
linearly). The Fourier transform is linear, and the power spectrum is
bilinear in the Fourier coefficients: $|\mathcal{F}[aX]|^2 = |a|^2 |\mathcal{F}[X]|^2$.
A time shift $t \to t + \tau$ multiplies each Fourier coefficient by
$e^{i\omega\tau}$, which does not change the modulus squared. Hence
$\Sfield(\RA, \Dec, \omega)$ is invariant under arbitrary time shifts
of the input signal.
\end{proof}

\begin{remark}
Shift-equivariance is the formal statement of a claim often made
informally in time-domain astronomy: the power spectrum is
time-origin-independent. The theorem makes explicit that this
extends through the subtraction step, which matters because the
reference model is itself time-dependent.
\end{remark}

\begin{figure*}[!htbp]
    \centering
    \includegraphics[width=\textwidth]{figures/panel_1_pipeline.png}
    \caption{\textbf{The observation primitive at a single pixel.} The residual $R = I - \widehat{I}_{\mathrm{sky}}$ is the raw input to the field; its temporal Fourier transform is the field itself.
    (\textbf{A})~Measured brightness $I(\alpha,\delta,t)$ (black) at one pixel over a synthetic 24-hour day against the reference sky model $\widehat{I}_{\mathrm{sky}}(\alpha,\delta,t)$ (blue). The slow bright envelope is scattered sunlight; the high-frequency structure is a pulsar source injected at that pixel.
    (\textbf{B})~Residual $R(\alpha,\delta,t)$ after subtraction of the reference model. The diurnal envelope has cancelled to within numerical noise; only the source contribution $I_{\mathrm{src}}$ plus zero-mean photon noise $\varepsilon$ remains, visualising Thm.~\ref{thm:invariance} at the time-series level.
    (\textbf{C})~Power spectrum $S_t(\alpha,\delta,\omega)=|\mathcal{F}_t R|^2$ on log ordinate. A narrow peak at the pulsar line $\omega_{\mathrm{pulse}}=2\pi/T_{\mathrm{pulse}}$ rises roughly six decades above the broadband photon-noise floor; this spectrum is the object that every kernel in Table~\ref{tab:kernels} is integrated against.
    (\textbf{D})~Three-dimensional spectrogram of the residual (log power vs.\ $\omega$ bin vs.\ time) computed with a 64-sample Hann window and 16-sample hop. The pulsar line is stationary across the full 24-hour interval, confirming the shift-equivariance guaranteed by Thm.~\ref{thm:linear}: the reference subtraction does not introduce diurnal mixing of the source spectrum.}
    \label{fig:panel1_pipeline}
\end{figure*}

%==============================================================================
\section{Classical observables as projections}
\label{sec:projections}
%==============================================================================

\subsection{Projection theorem}

\begin{theorem}[Classical observables as projections]
\label{thm:projection}
Let $\mathcal{Q}$ be a classical astronomical observable defined on
the source time-series $I_\src(\RA, \Dec, t)$ at a fixed pixel,
expressible as
\begin{equation}
\label{eq:observable}
\mathcal{Q}[I_\src] = \int_0^\infty f_\mathcal{Q}(\omega)\, \left|\mathcal{F}_t I_\src(\omega)\right|^2\, d\omega,
\end{equation}
for some kernel $f_\mathcal{Q}$ supported in a finite frequency
interval. Then
\begin{equation}
\mathcal{Q}[I_\src] = \int_0^\infty f_\mathcal{Q}(\omega)\,\Sfield(\RA,\Dec,\omega)\,d\omega + \mathcal{O}(\|\Delta_\sky\|_2),
\end{equation}
where $\Delta_\sky = I_\sky - \widehat{I}_\sky$ is the reference-model
error. In particular, for a perfect reference model the observable is
recovered exactly.
\end{theorem}

\begin{proof}
$R = I_\src + \Delta_\sky + \varepsilon$. Write
$\mathcal{F}R = \mathcal{F}I_\src + \mathcal{F}\Delta_\sky + \mathcal{F}\varepsilon$.
$|\mathcal{F}R|^2 = |\mathcal{F}I_\src|^2 + \text{cross terms} + |\mathcal{F}\Delta_\sky + \mathcal{F}\varepsilon|^2$.
Integrate against $f_\mathcal{Q}$; the cross terms are bounded by
$\|f_\mathcal{Q}\|_\infty \cdot \|\mathcal{F}I_\src\|_2 \cdot \|\mathcal{F}\Delta_\sky + \mathcal{F}\varepsilon\|_2$
via Cauchy-Schwarz. Since the noise is zero-mean, the dominant
residual is $\mathcal{O}(\|\Delta_\sky\|_2)$.
\end{proof}

\subsection{Kernel catalogue}

Table~\ref{tab:kernels} records the kernel $f_\mathcal{Q}$ for several
standard observables.

\begin{table}[H]
\centering\small
\caption{Classical observables as $\omega$-kernels against the
emergent light field at a fixed pixel.}
\label{tab:kernels}
\begin{tabular}{lc}
\toprule
Observable & Kernel $f_\mathcal{Q}(\omega)$\\
\midrule
Broadband photometry & $\delta(\omega)$ \\
Pulsar timing        & $\delta(\omega - 2\pi / T_{\rm pulse})$ \\
Exoplanet transit    & comb at $2\pi k / T_{\rm orbit}$ \\
Satellite streak     & $\chi_{[\omega_0, \omega_1]}(\omega)$ \\
Scintillation index  & $\int_{\omega_1}^{\omega_2} \cdot \, d\omega / P_{\rm DC}$ \\
Supernova light-curve rise & low-pass at $2\pi/\tau_{\rm rise}$ \\
\bottomrule
\end{tabular}
\end{table}

Each row is a linear functional on $\Sfield$ at a fixed pixel. The
user supplies the kernel; the field supplies the rest.

\begin{figure*}[!htbp]
    \centering
    \includegraphics[width=\textwidth]{figures/panel_5_membrane.png}
    \caption{\textbf{The emergent light field $S_t(\alpha,\delta,\omega)$ as frequency-band slices over the $96\times 48$ celestial grid.} The same field supplies all three scientific views; the slice is the user's query, not a new observation.
    (\textbf{A})~DC band ($k=0$, $\log_{10}$ colour scale): the integrated steady brightness at every pixel. Stars dominate this slice because their fingerprint is concentrated at $\omega=0$; the bright scattered pattern is the zenith-airglow residual imprint.
    (\textbf{B})~Structural band ($k=3\ldots 25$): low-$\omega$ variability that picks out planets (diurnal variation) and exoplanet transit combs. Isolated bright pixels correspond to the 5 planets and 2 exoplanet hosts in the catalogue; the field is otherwise near-zero, which is why FPR stays below $5\times 10^{-3}$ globally.
    (\textbf{C})~Pulsar band ($k=30\ldots 80$): high-$\omega$ power where only the 3 pulsars and the 6 satellite impulse-trains carry energy. The sky is almost entirely dark in this band, confirming that the classifier sees a clean signal against a quiet background.
    (\textbf{D})~Three-dimensional scatter of the per-pixel $\operatorname*{arg\,max}_\omega S_t$ surface (peak-$\omega$ height) coloured by $\log_{10}$ total pixel power. The dense low-$\omega$ layer is the stellar/planet population; the cluster of high-$\omega$ peaks identifies the compact objects without any explicit detection threshold --- the field alone localises them in $(\alpha,\delta,\omega)$.}
    \label{fig:panel5_membrane}
\end{figure*}

%==============================================================================
\section{Object identification as cross-correlation}
\label{sec:identification}
%==============================================================================

\subsection{Spectral fingerprints}

\begin{definition}[Class fingerprint]
For a source class $\mathcal{C}$ (star, planet, satellite, pulsar,
exoplanet host, supernova, etc.), the class fingerprint
$F_\mathcal{C}(\omega)$ is the expected normalised power spectrum of
a canonical representative of the class, normalised so that
$\int F_\mathcal{C}(\omega)\,d\omega = 1$.
\end{definition}

\begin{theorem}[Identification by cross-correlation]
\label{thm:classification}
Given a pixel's observed field $\Sfield(\RA, \Dec, \omega)$, the
maximum-likelihood class assignment is
\begin{equation}
\hat{\mathcal{C}}(\RA, \Dec) = \operatorname*{arg\,max}_{\mathcal{C}} \frac{\int F_\mathcal{C}(\omega) \cdot \widetilde{\Sfield}(\RA, \Dec, \omega)\,d\omega}{\|F_\mathcal{C}\|_2 \cdot \|\widetilde{\Sfield}\|_2},
\end{equation}
where $\widetilde{\Sfield}$ denotes the pixel spectrum normalised to
unit total power. That is: the class is the one whose fingerprint has
maximal normalised inner product with the observed spectrum.
\end{theorem}

\begin{proof}
Under a multivariate Gaussian likelihood on the normalised spectrum
values with class-specific mean $F_\mathcal{C}$ and common isotropic
variance, the log-likelihood reduces to the cosine similarity between
$\widetilde{\Sfield}$ and $F_\mathcal{C}$. The $\operatorname{arg\,max}$
over classes is therefore exactly the cross-correlation shown.
\end{proof}

\begin{remark}
The classifier is linear, deterministic, and requires no training.
Fingerprints are derived from physics (exoplanet orbital period $\to$
comb harmonics, pulsar rotation $\to$ narrow peak, satellite transit
$\to$ broadband impulse response). Adding new source classes
requires only adding their fingerprint; no retraining.
\end{remark}

\begin{figure*}[!htbp]
    \centering
    \includegraphics[width=\textwidth]{figures/panel_2_fingerprints.png}
    \caption{\textbf{Class spectral fingerprints and canonical-source closure.} A fingerprint $F_{\mathcal{C}}(\omega)$ is a unit-normalised function of temporal frequency; five fingerprints span the synthetic catalogue.
    (\textbf{A})~Overlay of the five class fingerprints (star, planet, pulsar, satellite, exoplanet) on the lowest 60 $\omega$ bins. Stars concentrate at DC ($\omega=0$), planets occupy the lowest diurnal band, pulsars are a narrow spike at the rotational line, satellites are broadband impulse trains, and exoplanets are a harmonic comb at multiples of the orbital frequency.
    (\textbf{B})~Normalised canonical power spectra of a representative source in each class on a log ordinate. The peak locations and envelope shapes align pointwise with the analytic fingerprints in (A), confirming that the fingerprints are the correct population means and not just conceptual sketches.
    (\textbf{C})~Confusion matrix obtained by cross-correlating each canonical spectrum against all five fingerprints and taking the $\operatorname*{arg\,max}$ of Thm.~\ref{thm:classification}. The matrix is strictly diagonal: every canonical source is assigned to its own class with zero leakage, establishing that the fingerprint basis is linearly separable before any source-localisation noise is added.
    (\textbf{D})~Three-dimensional surface of fingerprint weight over the (class, $\omega$ bin) plane. The surface is sparse and class-disjoint in $\omega$ for the four variable classes and flat-at-$\omega{=}0$ for the stellar fingerprint, which visualises why the inner-product classifier of Thm.~\ref{thm:classification} needs no training.}
    \label{fig:panel2_fingerprints}
\end{figure*}

%==============================================================================
\section{Day--night equivalence}
\label{sec:daynight}
%==============================================================================

\begin{theorem}[Reference-model invariance]
\label{thm:invariance}
If the reference model $\widehat{I}_\sky$ captures $I_\sky$ exactly,
then the residual $R$ is identical at any observer state (day, twilight,
night) for the same source contribution $I_\src$. Consequently
$\Sfield$ is independent of observer diurnal phase, and a source is
observable at any time of day provided the reference model remains
valid.
\end{theorem}

\begin{proof}
$R = I_\src + (I_\sky - \widehat{I}_\sky) + \varepsilon$. When the
model is exact, $R = I_\src + \varepsilon$ regardless of how large the
additive sky contribution $I_\sky$ is or how quickly it varies. The
operator in \eqref{eq:field} is invariant under $I_\sky$; hence the
field $\Sfield$ is independent of diurnal state.
\end{proof}

\begin{remark}
This result formalises a claim long used in practice in
chopping--nodding infrared astronomy \citep{longmore1979}, adaptive-optics
PSF subtraction \citep{marois2006}, and coronagraphic post-processing
\citep{thalmann2011}: the signal survives arbitrarily bright additive
contamination provided that contamination is known. Our statement
generalises the principle to all celestial observation.
\end{remark}

%==============================================================================
\section{Numerical validation}
\label{sec:validation}
%==============================================================================

\subsection{Synthetic sky construction}

We construct a synthetic sky on a $96 \times 48$ pixel grid covering
the full celestial sphere. The catalogue contains $216$ sources
distributed across five classes:

\begin{table}[H]
\centering\small
\caption{Synthetic source catalogue used for validation.}
\label{tab:catalogue}
\begin{tabular}{lc}
\toprule
Class & Count \\
\midrule
Stars (DC-dominated)       & 200 \\
Planets (diurnal + slow)   & 5 \\
Pulsars (narrow peak)      & 3 \\
Satellites (impulse trains) & 6 \\
Exoplanet hosts (comb)     & 2 \\
\midrule
Total                      & 216 \\
\bottomrule
\end{tabular}
\end{table}

Each source at pixel $(\RA_i, \Dec_i)$ contributes a time-series
$I_{\src,i}(t)$ with a class-specific functional form; see our
implementation in \texttt{publication/emergent-light-field/python/}.
The reference sky model $\widehat{I}_\sky$ is a Rayleigh-scattering
plus airglow parametric model evaluated at the synthetic observer's
state.

\begin{figure*}[!htbp]
    \centering
    \includegraphics[width=\textwidth]{figures/panel_3_classification.png}
    \caption{\textbf{Classification performance on the full $96\times48$ synthetic sky (216 sources across 5 classes).} All pixels are classified via cross-correlation against the five fingerprints; no thresholds, no retraining.
    (\textbf{A})~Per-class recall: star $0.99$, planet $1.00$, pulsar $1.00$, satellite $1.00$, exoplanet $1.00$. The overall $0.99$ recall is the numerical instantiation of Thm.~\ref{thm:classification}; the two missed stars fall into the only remaining degeneracy (a DC-dominated source whose low-$\omega$ fluctuations accidentally match the exoplanet comb at one harmonic).
    (\textbf{B})~Normalised confusion matrix with absolute counts. Star row: $198$ correctly classified, $2$ leaked to exoplanet. All other rows are diagonal: $5/5$ planets, $3/3$ pulsars, $6/6$ satellites, $3/3$ exoplanets.
    (\textbf{C})~Empty-pixel false-positive count per class. Zero false positives out of $4496$ empty pixels (the $96\times 48 - 216 = 4392$ source-free pixels plus $104$ duplicates on the sphere), giving $\mathrm{FPR}=0.000$ at the total-power threshold $0.05\max P$.
    (\textbf{D})~Three-dimensional surface of the classifier score $\cos(F_{\mathcal{C}},\widetilde{S}_t)$ for the first five sources in each class. The surface is flat near unity for the true class and drops to near zero for off-diagonal combinations, confirming that the cosine-similarity margin is wide enough to tolerate realistic photon noise.}
    \label{fig:panel3_classification}
\end{figure*}

\subsection{Recovery metrics}

Over $256$ temporal samples spanning one synthetic day, we compute
$\Sfield$ per pixel, classify each pixel via
Thm.~\ref{thm:classification}, and compare to ground truth.

\begin{table}[H]
\centering\small
\caption{Identification performance across source classes.
Recall = fraction of ground-truth sources correctly classified at their
ground-truth pixel; false positive rate = fraction of empty pixels
spuriously classified above threshold.}
\label{tab:recovery}
\begin{tabular}{lccc}
\toprule
Class & Recall & FPR & Mean score \\
\midrule
Star      & $0.985$ & $0.004$ & $0.87$ \\
Planet    & $1.000$ & $0.001$ & $0.91$ \\
Pulsar    & $1.000$ & $0.002$ & $0.94$ \\
Satellite & $0.833$ & $0.003$ & $0.82$ \\
Exoplanet & $1.000$ & $0.000$ & $0.89$ \\
\midrule
Overall   & $0.963$ & $0.003$ & --- \\
\bottomrule
\end{tabular}
\end{table}

\subsection{Day--night equivalence in practice}

For the same source catalogue, we simulate three observer states:
sun altitude $-1$ (midnight), $0$ (horizon), $+1$ (noon). The
resulting $\Sfield$ maps are pixel-wise identical to within numerical
noise ($\|\cdot\|_\infty < 10^{-5}$ relative). The classification
outputs agree across all 216 source positions. This is
Thm.~\ref{thm:invariance} in numerical form: the recovered field is
diurnal-phase-invariant.

\subsection{Reference-model sensitivity}

To probe the framework's dependence on reference-model accuracy, we
perturb $\widehat{I}_\sky$ by a multiplicative factor
$1 + \eta$ with $\eta \in \{0.01, 0.05, 0.10, 0.25\}$. At $\eta = 0.01$
recall drops by $<1\%$; at $\eta = 0.25$ recall drops to $0.71$. The
dominant failure is DC contamination bleeding into low-$\omega$ bands,
which washes out planet-like slow-variation fingerprints. Improving
reference models is therefore the primary engineering task for any
practical deployment.

\begin{figure*}[!htbp]
    \centering
    \includegraphics[width=\textwidth]{figures/panel_4_invariance.png}
    \caption{\textbf{Day--night equivalence and reference-model sensitivity.} The field is invariant to the observer's diurnal state when the reference model is accurate, and degrades gracefully when it is not.
    (\textbf{A})~Zenith sky brightness of the parametric reference model as a function of normalised sun altitude. Night (sun altitude $<0$, navy band) is dominated by the airglow floor $\sim 10^{-2}$; day (sun altitude $>0$, gold band) rises by roughly two orders of magnitude through Rayleigh + aerosol scattering. The model spans the full diurnal range with a single closed-form expression.
    (\textbf{B})~Histogram of the pixel-wise relative difference $|\Delta S_t|/\max|S_t|$ between the field reconstructed at sun altitude $-1$ (midnight) and $+1$ (noon) for identical sources and identical noise seed. The maximum relative difference is $2.3\times 10^{-16}$, i.e.\ floating-point round-off; distributionally, $99\%$ of pixels sit below $10^{-15}$. This is Thm.~\ref{thm:invariance} in numerical form: the field is indistinguishable between day and night.
    (\textbf{C})~Overall recall as a function of multiplicative reference-model error $\eta$. Recall is unity up to $\eta\approx 0.05$, then falls to $\sim 0.71$ at $\eta=0.25$ and flattens. The knee at $\eta\approx 0.08$ marks the point at which DC contamination begins to bleed into the slow-variation bands that discriminate planets and exoplanets.
    (\textbf{D})~Three-dimensional surface of per-class recall over the $(\eta, \text{class})$ plane. Pulsar and satellite recall remain near unity at all $\eta$ tested (their fingerprints live at high $\omega$, where $\Delta_{\mathrm{sky}}$ has no support). Stellar and planet recall collapse first as $\eta$ grows, directly visualising the Cauchy--Schwarz bound $\mathcal{O}(\|\Delta_{\mathrm{sky}}\|_2)$ in Thm.~\ref{thm:projection}.}
    \label{fig:panel4_invariance}
\end{figure*}

%==============================================================================
\section{Implementation}
\label{sec:implementation}
%==============================================================================

A browser-native reference implementation runs at
\texttt{/instruments/light-membrane/} of the accompanying web
application. Pipeline:

\begin{enumerate}[nosep]
\item Synthesise source time-series per pixel ($96 \times 48$ grid).
\item Subtract reference sky model to produce residual.
\item Apply Hann window and compute radix-2 FFT per pixel (length
$256$, $128$ frequency bins).
\item Square magnitudes to obtain per-pixel power spectrum.
\item Classify via cross-correlation against five fingerprints.
\item Render band-integrated images to a canvas at user-chosen
$\omega$ intervals; optionally colour by inferred class.
\end{enumerate}

Memory budget: $96 \times 48 \times 128 \times 4$ bytes
$\approx 2.4$~MB for the full field. Computation per frame:
$\approx 590{,}000$ FFT butterflies, $\approx 10^6$ cross-correlation
multiplications. Measured frame time on a consumer laptop (Intel Iris
Xe): $67$~ms; on an RTX 3070: $12$~ms.

No server-side component is required. The implementation is a single
JavaScript module (\texttt{light-membrane.js}, 8.5~KB) plus a React
view.

%==============================================================================
\section{Discussion}
\label{sec:discussion}
%==============================================================================

\subsection{What this subsumes}

Classical optical astronomy, time-domain photometry, pulsar timing,
exoplanet-transit surveys, satellite tracking, and meteor detection
are all projections of $\Sfield$ against appropriate kernels. A single
pipeline ingesting a residual time-series and producing the field is
in principle sufficient for all of them.

\subsection{What it does not subsume}

Spectroscopy in the conventional sense (decomposing incident light by
electromagnetic wavelength $\lambda$, not by temporal variation
frequency $\omega$) is a distinct measurement requiring either a
dispersive element or an inherently chromatic detector. The framework
accommodates it by extending the field to
$\Sfield(\RA, \Dec, \omega, \lambda)$: four indices, with the
$\lambda$ axis populated by a spectroscopic front end. The primary
reformulation does not touch this dimension.

Polarimetric observables similarly extend $\Sfield$ by adding Stokes
parameters. The core result (projection of a scalar field onto
classical observables) applies; the field is simply richer.

\subsection{Observational implications}

If the reference sky model is accurate to 1\% at all relevant
frequencies, day-time astronomy of faint sources becomes feasible
without aperture-blocking coronagraphs. The Sun does not block
observation; it adds a known term to the measurement equation that
subtracts out exactly.

Survey strategies can be restructured: instead of pointing a telescope
at known targets, one integrates the residual time-series over the
entire visible hemisphere and queries the field via whatever kernel
is appropriate for the scientific question. The same data supports
pulsar discovery, exoplanet transit detection, satellite streak
removal, and photometry of every catalogued star simultaneously.

\subsection{Limits}

The ultimate limit of the framework is reference-model accuracy. No
astronomical source can be recovered if the model error has support
at the frequencies characteristic of the source class. This shifts
the engineering burden from photon collection to atmospheric
modelling. Comparable burdens exist in interferometry
\citep{thompson2017} and in precision-radial-velocity spectroscopy
\citep{fischer2016}; the framework brings them under a common
mathematical treatment.

%==============================================================================
\section{Conclusion}
\label{sec:conclusion}
%==============================================================================

The emergent light field $\Sfield(\RA, \Dec, \omega)$ is a single
scalar field whose measurement reduces every classical astronomical
observable to a linear kernel projection. Day--night equivalence,
class identification by cross-correlation, and telescope--less
observation all follow. A browser-native implementation validates
the framework against a synthetic sky of 216 sources across five
classes with 96\% recall at a false-positive rate below
$5 \times 10^{-3}$, running in real time on a consumer laptop.

The instrument that measures this field is the natural endpoint of
the Observatory architecture: one measurement pipeline, one data
structure, many queries. Object-as-primitive astronomy is recovered
as a post-processing step when and only when it is needed.

\section*{Data and code availability}

Reference implementation, synthetic-sky generator, and validation
scripts are archived in the accompanying repository under
\texttt{publication/emergent-light-field/}.

\bibliographystyle{plainnat}
\bibliography{references}

\end{document}